\section{引言}

\subsection{地月空间任务规划的计算需求}

随着 Artemis 计划、Gateway 空间站、嫦娥探月工程以及商业月球任务的推进，
地月空间（cislunar space）正成为航天任务设计的热点区域。与近地轨道不同，
地月空间的动力学由地球与月球的双引力主导：拉格朗日平动点附近的周期轨道
（Halo、NRHO、DRO 等）没有解析解，其设计依赖三体问题框架下的数值方法——
微分修正、多重打靶、参数延拓与不变流形理论；而轨道之间的转移设计则需要在
脉冲、低能量与低推力等多种推进模式之间权衡。这类计算具有三个共同特征：
\textbf{模型多}（圆型限制性三体问题 CR3BP、星历 $N$ 体、双圆四体问题
BCR4BP 等多保真度模型并存）、\textbf{迭代重}（一次星历修正往往涉及上万段
轨道积分）、\textbf{链条长}（初猜、修正、高保真传播、控制仿真环环相扣）。

\subsection{定位：LLM + Agent 时代的数值基座}

e2m2e（Earth to Moon, Moon to Earth）面向上述需求，定位是
\textbf{地月空间任务规划的算法工具集基座}。在“LLM + Agent”式自主任务
规划系统中，大语言模型负责理解任务意图、分解与编排子任务；e2m2e 负责数值
计算的那一半——构建地月空间动力学模型，生成周期轨道族，设计轨道之间的
转移路径，并把结果画出来供检查。为此，e2m2e 在传统的 Python API 与 CLI
之外，原生提供 MCP（Model Context Protocol）服务：13 个任务级工具通过
stdio 暴露给 LLM Agent，工具清单由 Facade 方法元数据派生，计算产物自动
归档、\texttt{record\_id} 跨工具链接，使得“用自然语言设计一条 NRHO 并做
100 次 Monte Carlo 轨道保持仿真”成为一条可执行的指令。

\subsection{与现有工具的关系}

\begin{itemize}
  \item \textbf{GMAT / STK}：通用任务分析工具，以图形界面和脚本驱动为主，
        验证数据充分但难以被 Agent 程序化调用；
  \item \textbf{Orekit / poliastro}：通用轨道力学库，覆盖近地到深空，
        但不以地月三体问题为中心，也无任务级封装；
  \item \textbf{e2m2e}：专注地月空间，把“选族—初猜—修正—高保真传播—
        转移—保持”整条设计链封装为任务级接口，数值内核用 Rust 实现，
        并以 MCP 原生对接 LLM Agent。
\end{itemize}

e2m2e 不追求替代通用工具，而是与其形成互补：传播精度对齐 GMAT 与 DFH 至
亚百米量级（见第~\ref{sec:verification}~节），坐标与时间链用 GMAT 参考数据
交叉验证，同时提供通用工具所不具备的周期轨道族生成与 Agent 原生接口。
第~\ref{sec:comparison}~节基于对 10 个同类/生态工具的代码级调研，给出
系统的对比分析。

\subsection{本报告的组织}

第~\ref{sec:architecture}~节介绍系统架构：五个架构模块与“Python 编排 +
Rust 数值内核”的分层运行时；第~\ref{sec:dynamics}~节介绍时空基准、动力学
模型与高保真力模型；第~\ref{sec:orbit}~节与第~\ref{sec:transfer}~节分别
介绍周期轨道设计与转移轨道设计的核心算法；第~\ref{sec:control}~节介绍轨道
控制与 Monte Carlo 仿真；第~\ref{sec:interfaces}~节介绍 Facade/CLI/MCP
接口与轨道库；第~\ref{sec:comparison}~节将 e2m2e 与 GMAT、STK、ATK、
Orekit、pykep 等 10 个工具对比定位；第~\ref{sec:verification}~节给出
验证体系并总结展望。
